\SourceLeaf{535}\par Die von den Rechnern bis jetzt angewandten Arten zur Bestimmung von $\frac{\pi}{4}=(1)$ stellt \textsc{Gauss} in der folgenden Uebersicht zusammen
\begingroup\fontsize{10}{13}\selectfont\setlength{\tabcolsep}{3pt}
\begin{longtable}{@{}l p{.74\linewidth}@{}}
\textsc{Machin} & $(1)=4(5)-(239)$\par auch \textsc{Clausen} \\[3pt]
\textsc{Euler} & ${}=(2)+(3)$\par (\textsc{Euler} à \textsc{Goldbach} 1746 Mai 28) \\[3pt]
\textsc{Vega} & ${}=5(7)+2(\frac{79}{3})$\par (\textsc{Vega} Thesaurus logar. p. 633) \\[3pt]
\textsc{Vega} & ${}=2(3)+(7)$\par auch \textsc{Clausen} (Astr. Nachr. B. 25. S. 209) \\[3pt]
\textsc{Rutherford} & ${}=4(5)-(70)+(99)$\par (Philos. Trans. 1841. p. 283) \\[3pt]
\textsc{Dase} & ${}=(2)+(5)+(8)$\par (\textsc{Crelle} Journal. B. 27. S. 198) \\[3pt]
\textsc{Gauss}. 1. & ${}=12(18)+8(57)-5(239)$ \\[3pt]
\textsc{Gauss}. 2. & ${}=12(38)+20(57)+7(239)+24(268)$ \\[3pt]
\end{longtable}\endgroup
\par Die ersten Rechnungen für die Tafeln gehören der Zeit der Ausarbeitung der \textit{Disquiss. Arr.} an, sie sind dann besonders in den Jahren 1846 und 47 gefördert. Am 21. Juli 1847 waren 2283 Zerlegungen nach der hier wiedergegebenen Ordnung in Tafeln gebracht, die übrigen 169 sind später berechnet, und ich habe sie diesem Abdruck (der sich vom Original in der Einrichtung nur durch die des leichtern Satzes wegen statt der Potenzen angewandte Schreibweise der Wiederholung der Factoren unterscheidet) mit eingeordnet.
\par Die Manuscripte mit diesen letzten Rechnungen scheinen die Resultate in der Form zu enthalten, wie sie unmittelbar gefunden wurden. Die Reihenfolge, in welcher dabei die Zahlen $a$ auftreten, lässt vermuthen, dass nur für die kleinern die Theiler von $aa+1$ u.\,s.\,f. aufgesucht wurden, und dass die grössern Zahlen sich aus diesen durch Anwendung besonderer Kunstgriffe ergeben haben. Aufgezeichnet ist aber nur folgende Regel: \textit{Aus drei Zahlen $a$, $2a-n$, $2a+n$ findet sich eine vierte}
\[
\frac{4a^3-(nn-3)a}{nn+1}
\]
\textit{Diese ist immer eine ganze Zahl für $n=0$ und $n=1$, sonst nur}
\[
\begin{aligned}
&\textit{für}\quad a\equiv0\quad\textit{und}\quad\equiv\pm\surd-1\quad mod(nn+1)\quad\textit{wenn }n\textit{ gerade}\\
&\textit{und für}\quad a\equiv0\quad\textit{und}\quad\equiv\pm\surd-1\quad mod\frac{nn+1}{2}\quad\textit{wenn }n\textit{ ungerade}
\end{aligned}
\]
\begin{center}
\begin{tabular}{lrrr}
\textit{Beispiele} & $a=253,$ & $n=6,$ & $1750507$\\
 & $a=294,$ & $n=11,$ & $832902$\\
 & $a=119,$ & $n=1,$ & $3370437$\\
 & $a=57,$ & $n=3,$ & $74043$\\
 & $a=123,$ & $n=9,$ & $90657$
\end{tabular}
\end{center}
\par Zu der vierten Zahl gehören nemlich keine andern Primtheiler als zu den ersten dreien und davon sind auch nur diejenigen ungeraden Primtheiler ausgeschlossen, welche der Zahl $n$ zugehören.
\par Die Tabelle ‘\textit{Quadratorum myrias prima}’ enthält in der Zeile der Überschrift die Tausende und Hunderte, in der ersten senkrechten Spalte die Zehner und Einer der Grundzahl ferner in der letzten
